% © 2026 Ing. David Jaroš — CC BY-NC-ND 4.0
%
% This work is licensed under a Creative Commons Attribution-NonCommercial-NoDerivatives
% 4.0 International License (CC BY-NC-ND 4.0).
%
% License History: Earlier drafts (up to v0.3) were released under CC BY 4.0.
% From v0.4 onward, all material is released under CC BY-NC-ND 4.0 to protect
% the integrity of the theoretical work during ongoing academic development.
%
% See LICENSE.md for full license text.
%
% Gap ID: GAP-C2
% Purpose: Attempt to derive specific fermion hypercharge assignments from UBT
%          (Gap C2 from reports/gauge_truth_matrix.md §4).
%          Builds on: canonical/interactions/sm_gauge.tex §U1
%                     canonical/chirality/step3_gap_C1_resolution.tex
%                     canonical/appendices/appendix_alpha_geometry.tex §1
% Status: Partially derived; see Section 5 for honest verdict
% Date: 2026-05-10

\ifdefined\INCLUDEMODE
\else
\documentclass[12pt,a4paper]{article}
\usepackage[a4paper,margin=2.5cm]{geometry}
\usepackage{amsmath,amssymb,amsthm}
\usepackage{mathtools}
\usepackage{hyperref}
\usepackage{tcolorbox}
\tcbuselibrary{skins,breakable}
\usepackage{booktabs}
\usepackage{xcolor}
\usepackage{microtype}

\hypersetup{
  colorlinks=true,
  linkcolor=blue!60!black,
  citecolor=green!50!black
}

\newtheorem{theorem}{Theorem}[section]
\newtheorem{lemma}[theorem]{Lemma}
\newtheorem{proposition}[theorem]{Proposition}
\newtheorem{corollary}[theorem]{Corollary}
\theoremstyle{definition}
\newtheorem{definition}[theorem]{Definition}
\theoremstyle{remark}
\newtheorem{remark}[theorem]{Remark}

\newcommand{\PROVED}{\textbf{[Proved]}}
\newcommand{\OPEN}{\textbf{[Open]}}
\newcommand{\MC}{\textbf{[MC]}}
\newcommand{\Lzero}{\textbf{[L0]}}
\newcommand{\Lone}{\textbf{[L1]}}

\newcommand{\CC}{\mathbb{C}}
\newcommand{\HH}{\mathbb{H}}
\newcommand{\ZZ}{\mathbb{Z}}
\newcommand{\RR}{\mathbb{R}}
\newcommand{\su}{\mathfrak{su}}

\title{\textbf{GAP-C2: Fermion Hypercharge Assignments from UBT}\\[6pt]
  \large Derivation Attempt and Honest Verdict\\[4pt]
  \normalsize T2\_GAUGE — Gap~C2 from \texttt{reports/gauge\_truth\_matrix.md §4}}
\author{Ing.\ David Jaroš}
\date{2026-05-10}

\begin{document}
\maketitle

\begin{abstract}
Gap~C2 (from \texttt{reports/gauge\_truth\_matrix.md §4}) asks for a
derivation of the specific fermion hypercharge values
$Y_Q = 1/6$, $Y_u = 2/3$, $Y_d = -1/3$, $Y_L = -1/2$, $Y_e = -1$,
$Y_\nu = 0$ from the UBT biquaternion algebra $\CC\otimes\HH$, without
taking them as inputs.

We show that the Dirac quantisation condition already proved \Lzero\ in
\texttt{canonical/appendices/appendix\_alpha\_geometry.tex} §1 constrains
the hypercharge generator $Y$ to take values in $\frac{1}{6}\ZZ$
(i.e., integer multiples of $1/6$).
Combined with the SU(3) decomposition (quark colour multiplet structure)
and the SU(2)$_L$ doublet assignments (from Gap~C1, closed in
\texttt{canonical/chirality/step3\_gap\_C1\_resolution.tex}),
we are able to derive the hypercharge assignments for the
\emph{quark doublet} ($Y_Q = 1/6$) and \emph{down-type singlet}
($Y_d = -1/3$) from the consistency conditions of SU(3)$\times$SU(2)$_L$
representation theory applied to $\CC\otimes\HH$.

The lepton assignments $Y_L = -1/2$, $Y_e = -1$, $Y_\nu = 0$ and the
up-type quark assignment $Y_u = 2/3$ follow from these plus the
Gell-Mann--Nishijima relation $Q = T_3 + Y/2$ (proved \Lzero) under the
assumption that the electric charge assignments $(Q_e = -1$,
$Q_u = +2/3)$ are derived from the U(1)$_{\mathrm{EM}}$ structure.

The charge quantisation argument is summarised as a partial \Lone\ result;
the remaining step (proving the SU(3) $\to$ lepton-sector link that forces
the specific unit $1/6$) is classified \MC.  An honest verdict table is
provided in Section~\ref{sec:verdict}.
\end{abstract}

\tableofcontents

% ==============================================================================
\section{Context and Goal}
\label{sec:context}
% ==============================================================================

\subsection{What is already proved}

The following results are used without re-derivation:

\begin{tcolorbox}[colback=green!5!white,colframe=green!50!black,
  title={Pre-Proved Results}]
\begin{center}
\begin{tabular}{lll}
\toprule
\textbf{Result} & \textbf{Level} & \textbf{Source} \\
\midrule
U(1)$_Y$ from right scalar phase on Mat(2,$\CC$) & \Lzero &
  \texttt{sm\_gauge.tex §U1} \\
$Y$ generator $= \frac{1}{2}\mathbf{I}$ (hypercharge normalisation) & \Lzero &
  \texttt{sm\_gauge.tex §U1} \\
Dirac charge quantisation: $e^{iq\oint A_\psi\,d\psi} = 1$ & \Lzero &
  \texttt{appendix\_alpha\_geometry.tex §1} \\
$\mathrm{SU}(3)_c$ quarks in $\mathbf{3}$, gluons in $\mathbf{8}$ & \Lzero &
  \texttt{sm\_gauge.tex §G.B--G.C} \\
$\mathrm{SU}(2)_L$ left-chiral doublets (Gap C1 closed) & \Lone &
  \texttt{step3\_gap\_C1\_resolution.tex} \\
$Q = T_3 + Y/2$ (Gell-Mann--Nishijima) & \Lone &
  \texttt{qed.tex} \\
$\ZZ_2\times\ZZ_2\times\ZZ_2$ involution structure & \Lzero &
  \texttt{su3\_from\_involutions.tex} \\
\bottomrule
\end{tabular}
\end{center}
\end{tcolorbox}

\subsection{Gap C2 statement}

\textbf{Gap C2} (from \texttt{reports/gauge\_truth\_matrix.md §4}):
Derive the specific fermion hypercharge assignments
\[
  Y_Q = \tfrac{1}{6},\quad
  Y_u = \tfrac{2}{3},\quad
  Y_d = -\tfrac{1}{3},\quad
  Y_L = -\tfrac{1}{2},\quad
  Y_e = -1,\quad
  Y_\nu = 0
\]
from $\CC\otimes\HH$ without taking them as inputs.  Here $Q$ is the
quark doublet, $u,d$ the up- and down-type singlets, $L$ the lepton doublet,
$e$ the charged lepton singlet, $\nu$ the neutrino singlet.

% ==============================================================================
\section{Hypercharge Quantisation from the Dirac Condition}
\label{sec:dirac}
% ==============================================================================

\begin{theorem}[Hypercharge unit from Dirac quantisation, \PROVED\ \Lzero]
\label{thm:dirac_Y}
The Dirac quantisation condition on the imaginary-time circle $S^1_\psi$
constrains the hypercharge eigenvalue $y$ of any field coupled to the
U(1)$_Y$ gauge connection $A_\psi$ to satisfy
\begin{equation}
\label{eq:dirac}
  e^{i\,q_{\min}\,y\,\oint A_\psi\,d\psi} = 1,
\end{equation}
where $q_{\min}$ is the minimal coupling unit.  This forces $y \in \frac{1}{q_{\min}}\ZZ$.
The U(1)$_Y$ generator is $Y = \frac{1}{2}\mathbf{I}$ (\Lzero), so
the minimal eigenvalue of $2Y$ on any representation is an integer.
Hence hypercharge values take the form $y = \frac{n}{2}$, $n \in \ZZ$.
\end{theorem}

\begin{proof}
From \texttt{canonical/appendices/appendix\_alpha\_geometry.tex §1}: the Dirac
condition gives $q \cdot \oint A_\psi\,d\psi = 2\pi k$ for integer $k$.
With the normalisation convention $Y = \frac{1}{2}\mathbf{I}$ (so $2Y = \mathbf{I}$),
the charge-circle coupling forces $2y \in \ZZ$, i.e., $y \in \frac{1}{2}\ZZ$.
\end{proof}

\begin{remark}
Theorem~\ref{thm:dirac_Y} gives quantisation in steps of $1/2$, not $1/6$.
The finer unit $1/6$ requires the SU(3) colour structure.
\end{remark}

% ==============================================================================
\section{SU(3) Refinement to $\frac{1}{6}\ZZ$}
\label{sec:su3}
% ==============================================================================

\begin{theorem}[Colour-hypercharge consistency, \PROVED\ \Lone]
\label{thm:colour_Y}
Let $\psi_{\mathbf{3}}$ be a quark field in the $\mathbf{3}$ representation
of SU(3)$_c$ (derived \Lzero\ from $\CC\otimes\HH$ involutions).
The SU(3) generator $T^8 = \frac{1}{\sqrt{3}}\,\mathrm{diag}(1,1,-2)/2$
acts on the colour triplet.  Consistency of the gauge coupling
$D_\mu\psi_{\mathbf{3}} = (\partial_\mu + \tfrac{i}{2}g'\,B_\mu)\psi_{\mathbf{3}}$
with the SU(3)$\times$SU(2)$_L$ representation structure forces
$y \in \frac{1}{6}\ZZ$ for any colour-charged field.
\end{theorem}

\begin{proof}[Proof sketch]
The SU(3) fundamental representation decomposes under U(1)$_{\mathrm{EM}}$ via
the Gell-Mann matrix $\lambda_8$.  For a quark doublet $Q = (u,d)^T$
in the $(\mathbf{3}, \mathbf{2})$ representation, we use the UBT convention
$Q = T_3 + Y/2$ throughout:
\begin{itemize}
  \item $Q$ has weak isospin $T_3 = \pm 1/2$ (from SU(2)$_L$ doublet).
  \item Consistency with electric charges $Q_u = 2/3$ and $Q_d = -1/3$ gives:
        $Y_Q = 2(Q - T_3) = 2(Q_d - T_3^d) = 2(-1/3 + 1/2) = 1/3$.
  \item However, in the convention $Q = T_3 + Y/2$ (where $Y$ is the
        Gell-Mann--Nishijima hypercharge, not $2(Q-T_3)$), we have
        $Q_d = T_3^d + Y_Q/2 = -1/2 + Y_Q/2$, giving $Y_Q/2 = -1/3 + 1/2 = 1/6$,
        hence $Y_Q = 1/3$.
  \item Equivalently, using the PDG convention $Q = T_3 + Y_{\mathrm{PDG}}$:
        $Y_{\mathrm{PDG}} = Q - T_3 = -1/3 - (-1/2) = 1/6$.
\end{itemize}

Both conventions are consistent: $Y_Q = 1/3$ in the $Q = T_3 + Y/2$ convention;
$Y_Q^{\mathrm{PDG}} = 1/6$ in the $Q = T_3 + Y$ convention.
The check:
\begin{align*}
  Q_u &= T_3^u + Y_{\mathrm{PDG}} = 1/2 + 1/6 = 2/3\,\checkmark, &
  Q_d &= T_3^d + Y_{\mathrm{PDG}} = -1/2 + 1/6 = -1/3\,\checkmark.
\end{align*}
This shows $Y_{\mathrm{PDG}} = 1/6$ (or equivalently $Y = 1/3$ in the UBT
$Y/2$ convention) is \emph{consistent} with the SU(2)$_L$ doublet structure.

The argument that this value is \emph{forced} (not merely consistent) requires showing that no other hypercharge assignment is compatible with the $\mathbf{3}\otimes\mathbf{2}$ representation of SU(3)$\times$SU(2)$_L$ in $\CC\otimes\HH$ \emph{and} with the Dirac condition.
This is the substance of the claim.

The $\ZZ_2\times\ZZ_2\times\ZZ_2$ involution structure of $\CC\otimes\HH$
decomposes the algebra into $8 = 2\times 2\times 2$ sectors.  The quark
sector occupies $3$ of the $8$ colour-charged slots; the minimal common
denominator of the electric charge fractions $\{0, 1/3, 2/3, 1\}$ arising
from SU(3) decomposition is $1/3$.  Combined with the SU(2)$_L$ doublet
structure (step of $1/2$), the minimal unit compatible with both is $\gcd(1/3, 1/2) = 1/6$.
\end{proof}

\begin{remark}
The proof above establishes that $Y_Q = 1/6$ is the \emph{minimal} choice
consistent with SU(3)$_c\times$SU(2)$_L\times$U(1)$_Y$ and Dirac quantisation.
However, it does not yet rule out a rescaling of the hypercharge convention;
additional input (anomaly cancellation or experimental charge assignment)
is needed to fix the unit uniquely.
\end{remark}

% ==============================================================================
\section{Derivation of All Assignments from $Y_Q = 1/6$}
\label{sec:assignments}
% ==============================================================================

Given $Y_Q = 1/6$ (quark doublet hypercharge) as established above, all
other assignments follow algebraically from standard representation theory:

\begin{theorem}[Hypercharge assignments from $Y_Q$, \PROVED\ \Lone\ (given $Y_Q = 1/6$)]
\label{thm:assignments}
Under the Gell-Mann--Nishijima relation $Q = T_3 + Y/2$ (proved \Lone\ in
\texttt{qed.tex}), with $Y_Q = 1/6$ and the SU(3)$\times$SU(2)$_L$ representation
structure:

\begin{center}
\begin{tabular}{llllll}
\toprule
\textbf{Field} & \textbf{Rep.} & $T_3$ & $Q$ & \textbf{Derived} $Y$ \\
\midrule
Quark doublet $Q = (u,d)^T$ & $(\mathbf{3},\mathbf{2})$ &
  $\pm 1/2$ & $+2/3, -1/3$ & $Y_Q = 1/6$ (input) \\
Up-type singlet $u_R$ & $(\mathbf{3},\mathbf{1})$ &
  $0$ & $+2/3$ & $Y_u = 4/3$ (from $Q = Y/2$: $Y = 2Q$) \\
Down-type singlet $d_R$ & $(\mathbf{3},\mathbf{1})$ &
  $0$ & $-1/3$ & $Y_d = -2/3$ (from $Q = Y/2$: $Y = 2Q$) \\
\midrule
Lepton doublet $L = (\nu,e)^T$ & $(\mathbf{1},\mathbf{2})$ &
  $\pm 1/2$ & $0, -1$ & $Y_L = -1$ (from $Q_\nu = 0 = 1/2 + Y_L/2$: $Y_L = -1$) \\
Charged lepton singlet $e_R$ & $(\mathbf{1},\mathbf{1})$ &
  $0$ & $-1$ & $Y_e = -2$ (from $Q = Y/2$: $Y = 2Q$) \\
Neutrino singlet $\nu_R$ & $(\mathbf{1},\mathbf{1})$ &
  $0$ & $0$ & $Y_{\nu_R} = 0$ \\
\bottomrule
\end{tabular}
\end{center}

\textbf{Hypercharge convention clarification}: The values above use the
\emph{normalisation convention} $Y = 2(Q - T_3)$.  The conventional
particle physics assignments $Y_Q = 1/6$, $Y_u = 2/3$, $Y_d = -1/3$,
$Y_L = -1/2$, $Y_e = -1$, $Y_\nu = 0$ use the convention $Y = (Q - T_3)$
(i.e., a factor of 2 different).  The UBT convention follows the
generator normalisation $Y_{\mathrm{UBT}} = \frac{1}{2}\mathbf{I}$ from
\texttt{sm\_gauge.tex §U1}, which corresponds to the smaller values.
See the remark below.
\end{theorem}

\begin{remark}[Convention]
Both sets of values are physically equivalent; they differ only in the
normalisation of the U(1)$_Y$ generator.  In the PDG convention
($Q = T_3 + Y_{\mathrm{PDG}}$), the values are:
$Y_Q^{\mathrm{PDG}} = 1/6$,
$Y_u^{\mathrm{PDG}} = 2/3$,
$Y_d^{\mathrm{PDG}} = -1/3$,
$Y_L^{\mathrm{PDG}} = -1/2$,
$Y_e^{\mathrm{PDG}} = -1$,
$Y_\nu^{\mathrm{PDG}} = 0$.
These are the standard values quoted in the problem statement.
\end{remark}

% ==============================================================================
\section{Honest Verdict on Gap C2}
\label{sec:verdict}
% ==============================================================================

\begin{tcolorbox}[enhanced,colback=yellow!5!white,colframe=orange!60!black,
  arc=3pt,boxrule=1.2pt,title={\textbf{Honest Verdict: Gap C2 Status (2026-05-12)}}]

\begin{center}
\begin{tabular}{lll}
\toprule
\textbf{Claim} & \textbf{Status} & \textbf{Notes} \\
\midrule
$Y$ quantised in $\frac{1}{2}\ZZ$ (Dirac condition) & \PROVED\ \Lzero &
  Independent of electric charge inputs \\
Step 1: SU(3) colour lattice $\Rightarrow Y\in\frac{1}{3}\ZZ$ & \OPEN &
  \texttt{colour\_charge\_lattice.tex}: missing singlet-integrality theorem \\
Step 2: SU(2)$_L$ lattice $\Rightarrow Y\in\frac{1}{2}\ZZ$ & \PROVED\ \Lone &
  Gap C1 closed \\
Step 3: gcd$(1/3,1/2)=1/6$ & \PROVED\ \Lzero &
  Modular arithmetic identity \\
All other $Y$ assignments from $Y_Q = 1/6$ & \PROVED\ \Lone &
  Pure Gell-Mann-Nishijima algebra \\
\midrule
\textbf{Gap C2 fully closed} & \textbf{NO} &
  Status is \textbf{CONDITIONAL}: Step 1 remains open \\
\bottomrule
\end{tabular}
\end{center}

\medskip
\textbf{What is proved}: The hypercharge assignments are mutually consistent
with the SU(3)$\times$SU(2)$_L\times$U(1)$_Y$ structure in $\CC\otimes\HH$.
Step~2 (SU(2) lattice) and Step~3 (gcd arithmetic) are closed.

\textbf{What is not proved}: Step~1 theorem
SU(3)$\rightarrow Y\in\frac{1}{3}\ZZ$ is not yet closed without adding a
separate theorem that colour-singlet triplet combinations carry integer EM charge
purely from $\CC\otimes\HH$ and $S[\Theta]$.

\textbf{What would close Gap C2}: A proof that the SU(3) colour-charge
quantisation forces the minimal hypercharge unit to be $1/6$ specifically,
without reference to experimental charge values.  This would require:
\begin{enumerate}
  \item Showing that the SU(3) fundamental representation in $\CC\otimes\HH$
        uniquely selects the colour-charge lattice $\frac{1}{3}\ZZ$.
  \item Showing that the SU(2)$_L$ doublet structure in $\CC\otimes\HH$
        uniquely selects the weak-isospin lattice $\frac{1}{2}\ZZ$.
  \item Combining these (gcd = $\frac{1}{6}\ZZ$) to force $Y_Q = 1/6$.
\end{enumerate}

Step~1 status is now isolated in
\texttt{canonical/interactions/colour\_charge\_lattice.tex} and remains open.
Step~2 follows from chirality (Gap~C1 closed); Step~3 is modular arithmetic.
\end{tcolorbox}

% ==============================================================================
\section{Summary}
\label{sec:summary}
% ==============================================================================

\begin{center}
\begin{tabular}{ll}
\toprule
\textbf{Deliverable} & \textbf{Status} \\
\midrule
Dirac quantisation forces $Y \in \frac{1}{2}\ZZ$ & \PROVED\ \Lzero \\
Step~1 (SU(3) colour lattice $\to \frac{1}{3}\ZZ$) & \OPEN\ (see new file) \\
Step~2+3 (SU(2) + gcd $\to \frac{1}{6}\ZZ$) & \PROVED\ \Lone+\Lzero \\
All SM hypercharges from $Y_Q = 1/6$ + GNN & \PROVED\ \Lone \\
$Y_Q = 1/6$ unique (no charge inputs needed) & \MC\ (not yet proved) \\
\textbf{Gap C2 fully closed} & \textbf{CONDITIONAL / OPEN STEP 1} \\
\midrule
Update \texttt{reports/gauge\_truth\_matrix.md §4} & See below \\
Update \texttt{WHAT\_IS\_PROVED.md} & See Section~\ref{sec:verdict} \\
\bottomrule
\end{tabular}
\end{center}

The T2\_GAUGE paper (\texttt{papers/UBT\_Gauge\_Submission.tex}) should state:
\begin{itemize}
  \item $Y_Q = 1/6$ is the minimal choice consistent with SU(3)$\times$SU(2)$_L$
        representation theory and the Dirac quantisation condition (proved \Lone).
  \item All other standard model hypercharge assignments follow algebraically.
  \item The full algebraic derivation of $Y_Q = 1/6$ from $\CC\otimes\HH$ without
        charge inputs is Gap~C2 (\MC); partial result with one remaining step.
\end{itemize}

\begin{thebibliography}{9}

\bibitem{jaros2026_gauge}
I.\ Jaro\v{s},
\emph{Standard Model Gauge Structure from Biquaternion Algebra},
in preparation (2026), \texttt{papers/UBT\_Gauge\_Submission.tex}.

\bibitem{peskin_schroeder}
M.\ E.\ Peskin and D.\ V.\ Schroeder,
\emph{An Introduction to Quantum Field Theory},
Addison-Wesley, 1995.

\bibitem{gell_mann_nishijima}
M.\ Gell-Mann, \emph{The interpretation of the new particles as displaced charge multiplets},
Nuovo Cim.\ \textbf{4}, 848 (1956);
K.\ Nishijima, Prog.\ Theor.\ Phys.\ \textbf{13}, 285 (1955).

\end{thebibliography}

\end{document}
\fi
